
\section{Baseline Analysis}
\label{section: experimental setup - baseline analysis}

For our comparison, we need a DL model that is a State-of-the-Art work in the field of mmWave radar HPE.
% which ensures that both systems have comparable prior literature to build upon.
Furthermore, the DL model has to evaluate the same data that IAmMuse evaluates, the point clouds of the mmWave radar.
The output of the system should be a skeletal estimation, which can be transformed into an arm angle estimation in the same way as the ground truth.
With these considerations, the MARS model~\cite{an2021mars}, a State-of-the-Art model in the field of mmWave HPE and one of the first systems to generate skeletal estimations out of sparse mmWave point clouds, is selected as the ideal comparison model.

The MARS paper and this thesis both use a Kinect 2 for the ground truth and use mmWave radar point clouds as input.
One notable difference in the setups is the fact that MARS uses an older radar model, the IWR1443, as opposed to the IWR6843, though their output should be comparable.
Furthermore, the design of the MARS model has no temporal context and estimates each point cloud frame independently; therefore, it may struggle with sparse point cloud frames.


% \ednote{Either explain this in more detail here, or don't explain it here and do this in chapter 5.}
As a deep learning model, MARS's performance is directly linked to its training data.
Therefore, we consider multiple versions of MARS in the evaluation, each trained on a different dataset.
For the training of each of the models, we used a batch size of 128, with 150 epochs, in accordance with the example code provided on the GitHub repository of MARS~\cite{mars2021github}.
The different MARS models we considered each had different training and testing data, as shown in \cref{table: MARS train test}.
For the evaluation, we considered five different models:
\begin{itemize}
    \item \textbf{MARS baseline}: 
        This model demonstrates the performance of the MARS model if it uses the same training and testing data as IAmMuse.
        Therefore, it trains on the \textit{calibration data} of a \textbf{specific recording}, and it tests the \textit{interpretation data} of \textbf{the same recording}.
        The performance of this model is discussed in \cref{section: evaluation - baseline analysis}.
    \item \textbf{MARS pre-trained}: 
        This model shows the ability of MARS to generalise over datasets.
        For this, it uses the weights provided by the original MARS paper~\cite{an2021mars}, so it is trained on the MARS dataset, while it tests the \textit{interpretation data} of all recordings.
        The performance of this model is discussed in \cref{section: evaluation - larger models}.
    \item \textbf{MARS full calib}: 
        This model's training set consists of the \textit{calibration data} of \textbf{all recordings}, increasing the size of the training set to roughly one fifth that of the MARS paper.
        This model is also tested on the \textit{interpretation data} of all recordings.
        The performance of this model is discussed in \cref{section: evaluation - larger models}.
    \item \textbf{MARS partial standard}: 
        This model shows the ability of MARS to predict standard recordings when trained on very similar data.
        To do this, the model is trained on two-thirds of the \textit{interpretation data} of the standard recordings, and tested on the remaining one-third of the \textit{interpretation data} of the standard recordings.
        In this manner, the testing and training data are as similar as possible, which should result in optimal performance.
        The performance of this model is discussed in \cref{section: evaluation - iammuse vs best mars}.
    \item \textbf{MARS standard to free}: 
        This model tests the capacity of MARS to generalise from the standardised move set to a more realistic scenario.
        This model is trained with the interpretation data of \textit{all standardised recordings}, and is tested on the interpretation data of \textit{all free-play recordings}.
        The performance of this model is discussed in \cref{section: evaluation - iammuse vs best mars}.
\end{itemize}

By considering five distinct versions of MARS, we aim to provide a fair comparison, ensuring we do not disadvantage MARS with a single suboptimal training set.
Furthermore, the consideration of multiple models allows various failure modes of the MARS model to be shown.
Some of these failure modes are unique to the MARS models, but others are more generally found in Deep Learning models as a whole.


% \ednote{Jack zou hier een stukje toevoegen over waarom we 5 modelled geburiken. "We did not want to give a disadvantage to MARS by the type of training we gave it, so we considered multiple options, furthermore, we are able to see the various failure modes in this way".}


\begin{table}[!htb]
    \centering
    
    \resizebox{\linewidth}{!}{%
        \begin{tabular}{|p{0.30\linewidth}|p{0.5\linewidth}|p{0.5\linewidth}|}
        \hline
             \textbf{Model Name} & \textbf{Training Data} & \textbf{Testing Data}  \\
             \hline\hline
             MARS baseline  
                & Calibration data of \textit{one specific recording}.
                & Interpretation data of \textit{the same specific recording}.
             \\ \hline
             MARS pre-trained 
                & MARS data set.
                & Interpretation data of all recordings.
             \\ \hline
             MARS full calib 
                & Calibration data of all recordings.
                & Interpretation data of all recordings.
             \\ \hline
             MARS partial standard 
                & Interpretation data of 2/3 standard recordings. 
                & Interpretation data of the other 1/3 standard recordings.
             \\ \hline
             MARS standard to free 
                & Interpretation data of all standard recordings.
                & Interpretation data of all free-play recordings.
             \\ \hline
        \end{tabular}
    }
    \caption{The training and testing data for the various MARS models considered}
    \label{table: MARS train test}
\end{table}

% The specific training set and testing set for each of the considered models can be found in \cref{table: MARS train test}.
% Reasoning for the specific data sets chosen for these models can be found in \cref{chapter: evaluation}, specifically:
% \textit{MARS baseline} will be discussed in \cref{section: evaluation - baseline analysis}; 
% \textit{MARS paper} and \textit{MARS full calib} will be discussed in \cref{section: evaluation - larger models};
% and \textit{MARS partial standard} and \textit{MARS standard to free} will be discussed in \cref{section: evaluation - iammuse vs best mars}.







%     -[[ Introduction ]]
% - Why do we want to compare
% - What are some of the criteria a comparison system must fulfil
%   - Simple
%   - "comparable"
%   - Easy to get running.
% - 
% IAmMuse should also be compared to a similar Deep Learning model to see how well this algorithmic approach holds up to these more conventional systems, for this type of problem.
% There are a few criteria that an interpretation model should fulfil to be considered for use in this thesis.
% % early model
% One important consideration is that the model should be a relatively early attempt at HPE interpretation through DL methods.
% Since IAmMuse is the first attempt at algorithmic data interpretation for HPE, a fair comparison would put it up against an early DL system.
% % Same data
% The DL system should be able to run with the data that this thesis already has, so it needs to be able to take a Kinect skeletal estimation as a ground truth, and to take the pointclouds collected for this thesis as training data.
% % Easy to get going
% Lastly, for practical reasons, it should be relatively easy to get working, due to the general time constraints on a master's thesis.
% 
% 
% %     -[[ MARS ]]
% % - How does MARS fulfill the criteria
% % - Short synopsis of MARS system (e.g. Skeletal estimation through deep learning)
% % - Small notice of "training" versions
% % - 
% With these considerations, the MARS system \cite{an2021mars} was chosen.
% This is a relatively early system when it comes to mmWave radar interpretation for HPE, making it a fair comparison against IAmMuse.
% MARS was also trained by using the pointcloud output of a previous model of mmWave Radar from TI (the IWR1443) to the one that was used for this thesis (IWR6843).
% Next to that, the ground truth used in the MARS paper was also collected with a Kinect 2.
% Lastly, this model has already been used in our group, which meant that it was relatively easy to get it up and running.
% The exact method by which MARS produces a skeletal estimation is beyond the scope of this thesis, and can be found in the paper itself, however, it's still important to go over some basics.
% The MARS model has to \textit{train} on a set of data, the predictions it will make are based on this training.
% The MARS system also only takes into account a single frame at a time, therefore, it has no consensus of temporality.
% \ednote{refer back to \textit{related works} section for a brief explanation of MARS' inner workings}
% %     -[[ MARS variations ]]
% % - What "MARS versions"/trainings will we use
% % - Why were these chosen
% % - What are the expectations for each one?
% % - 
% As was already mentioned, MARS needs to be trained on a dataset to work. 
% This thesis will compare IAmMuse with multiple different trained versions of the MARS model, in order to get a better view of how the two systems compare.
% Table \ref{figure: MARS train test} shows the 5 different models that will be trained, and specifies the training and testing data used in each.
% For each specific MARS model, a batch size of 128 was used, with 150 epochs, in accordance with the example code provided in the MARS github repository\cite{mars2021github}.
% % Specific models
% A more detailed discussion of the different MARS models considered in this paper will be given in \cref{chapter: evaluation}, specifically,
% \textit{MARS baseline} will be discussed in \cref{section: evaluation - baseline analysis}; 
% \textit{MARS paper} and \textit{MARS full calib} will be discussed in \cref{section: evaluation - larger models};
% and \textit{MARS partial standard} and \textit{MARS standard to free} will be discussed in \cref{section: evaluation - iammuse vs best mars}.
% 
% \begin{table}[!htb]
    % \centering
    % 
    % \resizebox{\linewidth}{!}{%
        % \begin{tabular}{|p{0.30\linewidth}|p{0.5\linewidth}|p{0.5\linewidth}|}
        % \hline
             % \textbf{Model Name} & \textbf{Training Data} & \textbf{Testing Data}  \\
             % \hline\hline
             % MARS baseline  
                % & Calibration data of \textit{one specific recording}.
                % & Usage data of \textit{the same specific recording}.
             % \\ \hline
             % MARS paper 
                % & MARS data set.
                % & Usage data of all recordings.
             % \\ \hline
             % MARS full calib 
                % & Calibration data of all recordings.
                % & Usage data of all recordings.
             % \\ \hline
             % MARS partial standard 
                % & Usage data of 2/3 standard recordings. 
                % & Usage data of the other 1/3 standard recordings.
             % \\ \hline
             % MARS standard to free 
                % & Usage data of all standard recordings.
                % & Usage data of all free-play recordings.
             % \\ \hline
        % \end{tabular}
    % }
    % \caption{The training and testing data for the various MARS models considered}
    % \label{figure: MARS train test}
    % 
% \end{table}


%     -[[ ]]
% - 
% - 
% - 

%     -[[ ]]
% - 
% - 
% - 
%     -[[ ]]


% Double separation line and spacing
% \vspace{2em}
% \hrule
% \vspace{0.2em}
% \hrule
% \vspace{2em}


% Compare with a DL model
% Bring up the point of comparison, 
% Mention that we need a "baseline" for comparison.
% A baseline should be a deep learning model, since all current SOA systems are DL.
% The baseline should be a simple/early model, as IAmMuse is also a first exploration
% Propose MARS as a good model for comparison.

% Discuss MARS, the strengths and weaknesses according to its paper (single frame only, no temporality).
% Describe that we will want to train it on our own data.
% Briefly discuss the different training methods (single calib, all calib, MARS paper, standard to free, maybe partial standard).

% For this, use a table with the different train test data sets.
% Also, argue for each of these options, why they'd be interesting to look at.
% Give a small "alleup" for evaluation.

% Argue why MARS was used as a comparative model

% Discuss the training method of MARS (epocs, etc)

% Briefly describe the different training methods which will be discussed.
















% -- [ OLD ] --

% Describe the simplest training (trained on calib data)

% Describe some other "more advantageous" trainings we could do for Mars (all calib data, the MARS model itself, a fully representative section of the data).



